\chapter{How to Use the Reader Build}
\label{chap:reader-guide}

The main chapters already carry stage ranges, result-anchor references, and
claim-status summaries at the top. That is enough structure for a readable pass
if the reader uses the document in the right order.

\section{Recommended reading paths}

\subsection*{Path A: collaborator orientation}

Start with the anchor summary in the next chapter, then read only the parts
that matter for the current problem. Each part opens with:
\begin{itemize}[leftmargin=2.2em]
  \item the stage range it covers;
  \item a source/status table;
  \item the result-anchor families it mainly serves;
  \item a claim-status paragraph describing the status ceiling of the block.
\end{itemize}

\subsection*{Path B: theorem or formula audit}

Use the reader build to locate the right chapter and stage range, then move to
the archive build for the detailed trail. The normal sequence is:
\begin{enumerate}[leftmargin=2.2em]
  \item identify the relevant result-anchor family;
  \item read the theorem-level statement in the corresponding part;
  \item note the stage range and local status language;
  \item open the archive build if the claim needs full provenance or appendix
  detail;
  \item use the repo artifacts only when script-level or note-level checking is
  required.
\end{enumerate}

\subsection*{Path C: future-paper citation support}

Future papers should normally cite theorem-block results, not raw stage files.
The reader build is the fastest way to find the right result family and the
right status language. The archive build remains the place to verify the full
audit trail behind that citation.

\section{Reading discipline}

Three rules keep the reader build honest:
\begin{enumerate}[leftmargin=2.2em]
  \item A compact theorem block is not a status upgrade. If a chapter says a
  packet is reduced, numerical, or open, later prose must preserve that.
  \item A stage range is provenance, not a citation target by itself. Cite the
  result family first, then use stage ranges when the derivation history matters.
  \item A readable build is not a proof compression layer. When the claim is
  load-bearing, the archive build and repo-level audits are still the decisive
  references.
\end{enumerate}
